\textbf{Source metadata: author (MetaInlines)}

Owner-directed SplitZero analytic continuation

\textbf{Source metadata: date (MetaInlines)}

13 September 2026

\textbf{Source metadata: subtitle (MetaInlines)}

A uniform polynomial-norm estimate; integration with the four-volume criterion

\textbf{Source metadata: title (MetaInlines)}

Arithmetic endpoint bounds on the original τ-base source

\subsection{1. Result and exact source boundary}\label{endpoint-arithmetic_original--result-and-exact-source-boundary}

The formalization session's endpoint criterion selects an actual member of the canonical arithmetic family. This continuation proves one of its previously unproved analytic inputs. For the \textbf{unchanged arithmetic measure}, a fixed finite packet ENDPOINTSOURCE147F4C5896AEDFEFA0ED54AFSLOT0END, and integers ENDPOINTSOURCE147F4C5896AEDFEFA0ED54AFSLOT1END, there is a finite constant ENDPOINTSOURCE147F4C5896AEDFEFA0ED54AFSLOT2END, independent of ENDPOINTSOURCE147F4C5896AEDFEFA0ED54AFSLOT3END and ENDPOINTSOURCE147F4C5896AEDFEFA0ED54AFSLOT4END, such that

ENDPOINTSOURCE147F4C5896AEDFEFA0ED54AFSLOT5END

The norms are the original monic polynomial norms for ENDPOINTSOURCE147F4C5896AEDFEFA0ED54AFSLOT6END; the total mass remains ENDPOINTSOURCE147F4C5896AEDFEFA0ED54AFSLOT7END. No root-location hypothesis is used in this bound. The proof retains a compact-source mass, a lower convolution constant, and both exponential moments. These constants are defined explicitly below; they have not been numerically enclosed.

For the packet consisting exactly of a hypothetical nonreal off-line quartet, let

ENDPOINTSOURCE147F4C5896AEDFEFA0ED54AFSLOT8END

The existing exterior theorem gives ENDPOINTSOURCE147F4C5896AEDFEFA0ED54AFSLOT9END at every admitted degree. Combining (1) with the other session's endpoint theorem gives

ENDPOINTSOURCE147F4C5896AEDFEFA0ED54AFSLOT10END

Consequently such a quartet forces

ENDPOINTSOURCE147F4C5896AEDFEFA0ED54AFSLOT11END

The upper bound for this four-volume quantity is not proved here. In particular, (1) does not establish RH. It completes the polynomial-norm input in (3) on a specific admissible window and leaves the source/relation-volume comparison with its exact maps. Section 9 expresses that remaining quantity as two positive block log-determinants of actual new relation layers.

The finite quotient, canonical source, residue action and rank-two control are inputs from the preceding analytic notes. The new proofs in Sections 3--7 use only their explicitly defined arithmetic measure. The endpoint theorem is the formalizer's written result, not attributed to this continuation. Its finite telescopes are reported kernel-checked; the analytic estimates below have written proofs, not new Lean certificates.

\subsection{2. The base, source, and both kinds of quotient}\label{endpoint-arithmetic_original--the-base-source-and-both-kinds-of-quotient}

Retain the original marked coefficient square

ENDPOINTSOURCE147F4C5896AEDFEFA0ED54AFSLOT12END

Here ENDPOINTSOURCE147F4C5896AEDFEFA0ED54AFSLOT13END, ENDPOINTSOURCE147F4C5896AEDFEFA0ED54AFSLOT14END, ENDPOINTSOURCE147F4C5896AEDFEFA0ED54AFSLOT15END, and ENDPOINTSOURCE147F4C5896AEDFEFA0ED54AFSLOT16END. The target of ENDPOINTSOURCE147F4C5896AEDFEFA0ED54AFSLOT17END is infinite. The construction remains over the original pointed base ENDPOINTSOURCE147F4C5896AEDFEFA0ED54AFSLOT18END; no new scalar zero is adjoined at a tensor degree or at a polynomial relation layer.

Keep ENDPOINTSOURCE147F4C5896AEDFEFA0ED54AFSLOT19END in degrees ENDPOINTSOURCE147F4C5896AEDFEFA0ED54AFSLOT20END, ENDPOINTSOURCE147F4C5896AEDFEFA0ED54AFSLOT21END, and ENDPOINTSOURCE147F4C5896AEDFEFA0ED54AFSLOT22END. The source spaces and the Fourier transform remain those of the earlier theta construction. In particular,

ENDPOINTSOURCE147F4C5896AEDFEFA0ED54AFSLOT23END

Let ENDPOINTSOURCE147F4C5896AEDFEFA0ED54AFSLOT24END be a fixed finite product of zeros of ENDPOINTSOURCE147F4C5896AEDFEFA0ED54AFSLOT25END, with the full order of each selected zero retained, and write ENDPOINTSOURCE147F4C5896AEDFEFA0ED54AFSLOT26END. Define

ENDPOINTSOURCE147F4C5896AEDFEFA0ED54AFSLOT27END

The quotient ENDPOINTSOURCE147F4C5896AEDFEFA0ED54AFSLOT28END is entire, including at the selected centres. It is not the meromorphic expression evaluated before removing its actual removable factors. For ENDPOINTSOURCE147F4C5896AEDFEFA0ED54AFSLOT29END, the same analytic definitions are meaningful, but the finite arithmetic packet is zero. Only the nonempty case is used for the finite Gram inverses below.

The source map is

ENDPOINTSOURCE147F4C5896AEDFEFA0ED54AFSLOT30END

The original cyclic module is ENDPOINTSOURCE147F4C5896AEDFEFA0ED54AFSLOT31END. Its arithmetic inclusion keeps the Taylor unit:

ENDPOINTSOURCE147F4C5896AEDFEFA0ED54AFSLOT32END

The supplied source identities are

ENDPOINTSOURCE147F4C5896AEDFEFA0ED54AFSLOT33END

The actual norm of (8) is

ENDPOINTSOURCE147F4C5896AEDFEFA0ED54AFSLOT34END

For ENDPOINTSOURCE147F4C5896AEDFEFA0ED54AFSLOT35END, the monic norm is the quotient norm of the leading-coefficient map

ENDPOINTSOURCE147F4C5896AEDFEFA0ED54AFSLOT36END

The arithmetic quotient is the other specified map from the same source:

ENDPOINTSOURCE147F4C5896AEDFEFA0ED54AFSLOT37END

for ENDPOINTSOURCE147F4C5896AEDFEFA0ED54AFSLOT38END, with the preceding relation space zero at ENDPOINTSOURCE147F4C5896AEDFEFA0ED54AFSLOT39END. Multiplication by the monic ENDPOINTSOURCE147F4C5896AEDFEFA0ED54AFSLOT40END carries ENDPOINTSOURCE147F4C5896AEDFEFA0ED54AFSLOT41END isomorphically to ENDPOINTSOURCE147F4C5896AEDFEFA0ED54AFSLOT42END, with leading coefficient unchanged. The map ENDPOINTSOURCE147F4C5896AEDFEFA0ED54AFSLOT43END kills that very relation. Thus (12), (13), and multiplication by ENDPOINTSOURCE147F4C5896AEDFEFA0ED54AFSLOT44END specify the relation between the degree-line norm and the arithmetic quotient metric. One norm is not substituted for the other.

At each admitted original support label ENDPOINTSOURCE147F4C5896AEDFEFA0ED54AFSLOT45END, the lifted quotient is ENDPOINTSOURCE147F4C5896AEDFEFA0ED54AFSLOT46END. An original relation has image ENDPOINTSOURCE147F4C5896AEDFEFA0ED54AFSLOT47END; external absence has image ENDPOINTSOURCE147F4C5896AEDFEFA0ED54AFSLOT48END. Its norm before this projection remains available through (11).

\subsection{3. A quantitative local arithmetic mass lemma}\label{endpoint-arithmetic_original--a-quantitative-local-arithmetic-mass-lemma}

\subsubsection{3.1 Uniform polynomial growth on a fixed strip}\label{endpoint-arithmetic_original--uniform-polynomial-growth-on-a-fixed-strip}

There is a constant ENDPOINTSOURCE147F4C5896AEDFEFA0ED54AFSLOT49END such that, for ENDPOINTSOURCE147F4C5896AEDFEFA0ED54AFSLOT50END, the function

ENDPOINTSOURCE147F4C5896AEDFEFA0ED54AFSLOT51END

is holomorphic on and inside the ellipse

ENDPOINTSOURCE147F4C5896AEDFEFA0ED54AFSLOT52END

and has modulus at most ENDPOINTSOURCE147F4C5896AEDFEFA0ED54AFSLOT53END there.

The ellipse has real and imaginary semiaxes ENDPOINTSOURCE147F4C5896AEDFEFA0ED54AFSLOT54END and ENDPOINTSOURCE147F4C5896AEDFEFA0ED54AFSLOT55END. Hence its image in the ENDPOINTSOURCE147F4C5896AEDFEFA0ED54AFSLOT56END-plane lies in ENDPOINTSOURCE147F4C5896AEDFEFA0ED54AFSLOT57END and ENDPOINTSOURCE147F4C5896AEDFEFA0ED54AFSLOT58END. The pole of ENDPOINTSOURCE147F4C5896AEDFEFA0ED54AFSLOT59END is outside this set when ENDPOINTSOURCE147F4C5896AEDFEFA0ED54AFSLOT60END.

For ENDPOINTSOURCE147F4C5896AEDFEFA0ED54AFSLOT61END, use the identity

ENDPOINTSOURCE147F4C5896AEDFEFA0ED54AFSLOT62END

It follows by summing the steps of ENDPOINTSOURCE147F4C5896AEDFEFA0ED54AFSLOT63END for ENDPOINTSOURCE147F4C5896AEDFEFA0ED54AFSLOT64END and then by analytic continuation for ENDPOINTSOURCE147F4C5896AEDFEFA0ED54AFSLOT65END, ENDPOINTSOURCE147F4C5896AEDFEFA0ED54AFSLOT66END. Its integral is bounded by ENDPOINTSOURCE147F4C5896AEDFEFA0ED54AFSLOT67END, giving ENDPOINTSOURCE147F4C5896AEDFEFA0ED54AFSLOT68END on the indicated part of the ellipse. For ENDPOINTSOURCE147F4C5896AEDFEFA0ED54AFSLOT69END, the functional equation

ENDPOINTSOURCE147F4C5896AEDFEFA0ED54AFSLOT70END

and the vertical-strip gamma estimate give ENDPOINTSOURCE147F4C5896AEDFEFA0ED54AFSLOT71END. On the stated strip this exponent is below ENDPOINTSOURCE147F4C5896AEDFEFA0ED54AFSLOT72END. The common power ENDPOINTSOURCE147F4C5896AEDFEFA0ED54AFSLOT73END in (14) is therefore a valid upper bound with one fixed constant. The gamma estimate and functional equation are the classical inputs recorded in DLMF §§5.11 and 25.4; none assumes RH.

Every disc of radius one about a point of ENDPOINTSOURCE147F4C5896AEDFEFA0ED54AFSLOT74END lies inside the ellipse. The maximum principle and Cauchy's derivative formula therefore also give

ENDPOINTSOURCE147F4C5896AEDFEFA0ED54AFSLOT75END

\subsubsection{3.2 Propagate a nonzero value to the critical-line interval}\label{endpoint-arithmetic_original--propagate-a-nonzero-value-to-the-critical-line-interval}

Put

ENDPOINTSOURCE147F4C5896AEDFEFA0ED54AFSLOT76END

The absolutely convergent reciprocal Dirichlet series gives ENDPOINTSOURCE147F4C5896AEDFEFA0ED54AFSLOT77END, so ENDPOINTSOURCE147F4C5896AEDFEFA0ED54AFSLOT78END. Let ENDPOINTSOURCE147F4C5896AEDFEFA0ED54AFSLOT79END.

Apply the maximum principle on ENDPOINTSOURCE147F4C5896AEDFEFA0ED54AFSLOT80END to ENDPOINTSOURCE147F4C5896AEDFEFA0ED54AFSLOT81END. Equivalently subtract from its logarithmic modulus the harmonic function linear in ENDPOINTSOURCE147F4C5896AEDFEFA0ED54AFSLOT82END with boundary values ENDPOINTSOURCE147F4C5896AEDFEFA0ED54AFSLOT83END and ENDPOINTSOURCE147F4C5896AEDFEFA0ED54AFSLOT84END. Zeros are handled by the usual subharmonic maximum principle. At ENDPOINTSOURCE147F4C5896AEDFEFA0ED54AFSLOT85END this gives

ENDPOINTSOURCE147F4C5896AEDFEFA0ED54AFSLOT86END

Here ENDPOINTSOURCE147F4C5896AEDFEFA0ED54AFSLOT87END and ENDPOINTSOURCE147F4C5896AEDFEFA0ED54AFSLOT88END. Consequently

ENDPOINTSOURCE147F4C5896AEDFEFA0ED54AFSLOT89END

This is the three-circles argument with every coordinate map and radius fixed. No polynomial approximation to ENDPOINTSOURCE147F4C5896AEDFEFA0ED54AFSLOT90END or assumption about its zeros is used.

Choose a point attaining ENDPOINTSOURCE147F4C5896AEDFEFA0ED54AFSLOT91END. By (18), on a one-sided subinterval of length ENDPOINTSOURCE147F4C5896AEDFEFA0ED54AFSLOT92END inside ENDPOINTSOURCE147F4C5896AEDFEFA0ED54AFSLOT93END, the modulus is at least ENDPOINTSOURCE147F4C5896AEDFEFA0ED54AFSLOT94END. Thus

ENDPOINTSOURCE147F4C5896AEDFEFA0ED54AFSLOT95END

For ENDPOINTSOURCE147F4C5896AEDFEFA0ED54AFSLOT96END, the integral in (22) is continuous and strictly positive: otherwise ENDPOINTSOURCE147F4C5896AEDFEFA0ED54AFSLOT97END would vanish on an interval and hence identically on its holomorphic domain. Its minimum on that compact interval is positive. Combining it with (22) proves the all-real statement

ENDPOINTSOURCE147F4C5896AEDFEFA0ED54AFSLOT98END

The exponent ENDPOINTSOURCE147F4C5896AEDFEFA0ED54AFSLOT99END is deliberately an explicit sufficient exponent, not an optimal local mean-square estimate. The proof is a classical analytic propagation argument specialized here to the needed arithmetic interval.

\subsection{4. A lower envelope for the actual convolution, after three factors}\label{endpoint-arithmetic_original--a-lower-envelope-for-the-actual-convolution-after-three-factors}

\subsubsection{\texorpdfstring{4.1 Local mass of the unchanged ENDPOINTSOURCE147F4C5896AEDFEFA0ED54AFSLOT100END}{4.1 Local mass of the unchanged }}\label{endpoint-arithmetic_original--local-mass-of-the-unchanged-gh}

The gamma estimate, with a positive lower constant valid also on compact sets, implies

ENDPOINTSOURCE147F4C5896AEDFEFA0ED54AFSLOT101END

Let ENDPOINTSOURCE147F4C5896AEDFEFA0ED54AFSLOT102END, so ENDPOINTSOURCE147F4C5896AEDFEFA0ED54AFSLOT103END. The entire identity ENDPOINTSOURCE147F4C5896AEDFEFA0ED54AFSLOT104END makes this bound valid at selected centres as well: ENDPOINTSOURCE147F4C5896AEDFEFA0ED54AFSLOT105END. The factors in (6), including ENDPOINTSOURCE147F4C5896AEDFEFA0ED54AFSLOT106END, then give a positive constant ENDPOINTSOURCE147F4C5896AEDFEFA0ED54AFSLOT107END with

ENDPOINTSOURCE147F4C5896AEDFEFA0ED54AFSLOT108END

Equation (25) is a weaker bound than the available factor ENDPOINTSOURCE147F4C5896AEDFEFA0ED54AFSLOT109END; it follows by an inequality, not by changing the density. On ENDPOINTSOURCE147F4C5896AEDFEFA0ED54AFSLOT110END, its exponential and polynomial weights are bounded below by ENDPOINTSOURCE147F4C5896AEDFEFA0ED54AFSLOT111END and ENDPOINTSOURCE147F4C5896AEDFEFA0ED54AFSLOT112END, respectively. Hence (23) proves

ENDPOINTSOURCE147F4C5896AEDFEFA0ED54AFSLOT113END

\subsubsection{4.2 The positive compact convolution factor}\label{endpoint-arithmetic_original--the-positive-compact-convolution-factor}

The function ENDPOINTSOURCE147F4C5896AEDFEFA0ED54AFSLOT114END is continuous, bounded, integrable, and positive outside a discrete set. Thus ENDPOINTSOURCE147F4C5896AEDFEFA0ED54AFSLOT115END is continuous and strictly positive at every real point. Indeed its integrand is positive almost everywhere; continuity follows from translation continuity in ENDPOINTSOURCE147F4C5896AEDFEFA0ED54AFSLOT116END and boundedness of the other factor. Define the \textbf{actual} compact quantities

ENDPOINTSOURCE147F4C5896AEDFEFA0ED54AFSLOT117END

Both are determined by the original density; neither is assigned the value one. Restricting a nonnegative convolution integral gives

ENDPOINTSOURCE147F4C5896AEDFEFA0ED54AFSLOT118END

The omitted part is the nonnegative integral over ENDPOINTSOURCE147F4C5896AEDFEFA0ED54AFSLOT119END and remains part of ENDPOINTSOURCE147F4C5896AEDFEFA0ED54AFSLOT120END. We use a lower bound, not a new compactly supported density.

For ENDPOINTSOURCE147F4C5896AEDFEFA0ED54AFSLOT121END, restrict the remaining ENDPOINTSOURCE147F4C5896AEDFEFA0ED54AFSLOT122END variables to ENDPOINTSOURCE147F4C5896AEDFEFA0ED54AFSLOT123END and use (28). Their sum has absolute value at most ENDPOINTSOURCE147F4C5896AEDFEFA0ED54AFSLOT124END. This gives the central analytic estimate

ENDPOINTSOURCE147F4C5896AEDFEFA0ED54AFSLOT125END

All mass and tensor factors are explicit. For ENDPOINTSOURCE147F4C5896AEDFEFA0ED54AFSLOT126END the factor ENDPOINTSOURCE147F4C5896AEDFEFA0ED54AFSLOT127END is the empty product. No statement for ENDPOINTSOURCE147F4C5896AEDFEFA0ED54AFSLOT128END is inferred by using a negative exponent. The amplification below only needs ENDPOINTSOURCE147F4C5896AEDFEFA0ED54AFSLOT129END.

\subsection{5. Upper moments and exact lower monic norms}\label{endpoint-arithmetic_original--upper-moments-and-exact-lower-monic-norms}

Fix ENDPOINTSOURCE147F4C5896AEDFEFA0ED54AFSLOT130END and retain both one-factor Laplace values

ENDPOINTSOURCE147F4C5896AEDFEFA0ED54AFSLOT131END

They are finite by the original gamma decay. Fubini and positivity give

ENDPOINTSOURCE147F4C5896AEDFEFA0ED54AFSLOT132END

The elementary exponential-series inequality ENDPOINTSOURCE147F4C5896AEDFEFA0ED54AFSLOT133END for ENDPOINTSOURCE147F4C5896AEDFEFA0ED54AFSLOT134END yields

ENDPOINTSOURCE147F4C5896AEDFEFA0ED54AFSLOT135END

The trial polynomial is ENDPOINTSOURCE147F4C5896AEDFEFA0ED54AFSLOT136END, whose value on ENDPOINTSOURCE147F4C5896AEDFEFA0ED54AFSLOT137END is ENDPOINTSOURCE147F4C5896AEDFEFA0ED54AFSLOT138END. Its phase has modulus one and remains in the source map; the generator ENDPOINTSOURCE147F4C5896AEDFEFA0ED54AFSLOT139END has not been shifted to define a different spectral problem.

For every leading coefficient of modulus one, the exact least squared norm of a degree-ENDPOINTSOURCE147F4C5896AEDFEFA0ED54AFSLOT140END polynomial on ENDPOINTSOURCE147F4C5896AEDFEFA0ED54AFSLOT141END with Lebesgue measure is

ENDPOINTSOURCE147F4C5896AEDFEFA0ED54AFSLOT142END

To check all constants, Rodrigues' formula gives the Legendre polynomial ENDPOINTSOURCE147F4C5896AEDFEFA0ED54AFSLOT143END with leading coefficient ENDPOINTSOURCE147F4C5896AEDFEFA0ED54AFSLOT144END and squared norm ENDPOINTSOURCE147F4C5896AEDFEFA0ED54AFSLOT145END. Its scaled monic version on ENDPOINTSOURCE147F4C5896AEDFEFA0ED54AFSLOT146END is the unique minimizer: subtract it from any competing monic polynomial, and the difference has lower degree and is orthogonal to it. For the source coordinate ENDPOINTSOURCE147F4C5896AEDFEFA0ED54AFSLOT147END, multiply this minimizer by the retained leading phase ENDPOINTSOURCE147F4C5896AEDFEFA0ED54AFSLOT148END. This proves (33) at the original type. See also DLMF §§18.3 and 18.5.

Restricting (11) to ENDPOINTSOURCE147F4C5896AEDFEFA0ED54AFSLOT149END and using (29) proves

ENDPOINTSOURCE147F4C5896AEDFEFA0ED54AFSLOT150END

The observation has the explicit isometric splitting into the inside and outside interval components. Addition recovers the full source function, and their squared norms add. In the split lift each zero component is the supported zero at its label. The outside component has not been declared absent.

\subsection{6. The norm-window estimate is proved}\label{endpoint-arithmetic_original--the-norm-window-estimate-is-proved}

Set ENDPOINTSOURCE147F4C5896AEDFEFA0ED54AFSLOT151END and ENDPOINTSOURCE147F4C5896AEDFEFA0ED54AFSLOT152END in (34), and use ENDPOINTSOURCE147F4C5896AEDFEFA0ED54AFSLOT153END in (32). The exact finite bound is

ENDPOINTSOURCE147F4C5896AEDFEFA0ED54AFSLOT154END

For ENDPOINTSOURCE147F4C5896AEDFEFA0ED54AFSLOT155END, this is at most ENDPOINTSOURCE147F4C5896AEDFEFA0ED54AFSLOT156END. Here is one explicit sufficient constant, with ENDPOINTSOURCE147F4C5896AEDFEFA0ED54AFSLOT157END:

ENDPOINTSOURCE147F4C5896AEDFEFA0ED54AFSLOT158END

Proof of the uniform estimate: ENDPOINTSOURCE147F4C5896AEDFEFA0ED54AFSLOT159END in (33) gives ENDPOINTSOURCE147F4C5896AEDFEFA0ED54AFSLOT160END. Also ENDPOINTSOURCE147F4C5896AEDFEFA0ED54AFSLOT161END, ENDPOINTSOURCE147F4C5896AEDFEFA0ED54AFSLOT162END, ENDPOINTSOURCE147F4C5896AEDFEFA0ED54AFSLOT163END, and ENDPOINTSOURCE147F4C5896AEDFEFA0ED54AFSLOT164END. After taking the ENDPOINTSOURCE147F4C5896AEDFEFA0ED54AFSLOT165END-th root the factorial and monic factors contribute at most ENDPOINTSOURCE147F4C5896AEDFEFA0ED54AFSLOT166END. The mass factors contribute at most the square root in (36); the exponential contributes at most ENDPOINTSOURCE147F4C5896AEDFEFA0ED54AFSLOT167END. Finally ENDPOINTSOURCE147F4C5896AEDFEFA0ED54AFSLOT168END gives ENDPOINTSOURCE147F4C5896AEDFEFA0ED54AFSLOT169END. This proves (1).

The constant is intentionally not optimized. Its dependence on the fixed packet is allowed; its independence of the polynomial and tensor degrees is what is needed. Defining it via positive compact minima does not claim to have provided a numerical interval enclosure for those minima.

For a general source window the same two inequalities give, without asymptotic notation,

ENDPOINTSOURCE147F4C5896AEDFEFA0ED54AFSLOT170END

\subsection{7. Compose the proved bound with the formalizer's endpoint theorem}\label{endpoint-arithmetic_original--compose-the-proved-bound-with-the-formalizers-endpoint-theorem}

For a nonempty cyclic packet ENDPOINTSOURCE147F4C5896AEDFEFA0ED54AFSLOT171END, let ENDPOINTSOURCE147F4C5896AEDFEFA0ED54AFSLOT172END be its original canonical quotient metric, ENDPOINTSOURCE147F4C5896AEDFEFA0ED54AFSLOT173END, and ENDPOINTSOURCE147F4C5896AEDFEFA0ED54AFSLOT174END the original rank-two control allowance. The new endpoint theorem supplied by PR \#23 is

ENDPOINTSOURCE147F4C5896AEDFEFA0ED54AFSLOT175END

Its proof multiplies the local Toda bounds, applies concavity of ENDPOINTSOURCE147F4C5896AEDFEFA0ED54AFSLOT176END, and uses the two exact telescopes. The zero-contraction case is handled before taking logarithms. This result is credited to the parallel formalization/research session. The proof and the distinction between checked finite components and written Jensen assembly are retained in its source.

Choose ENDPOINTSOURCE147F4C5896AEDFEFA0ED54AFSLOT177END and ENDPOINTSOURCE147F4C5896AEDFEFA0ED54AFSLOT178END. Equations (1) and (38) yield the actual canonical selection estimate

ENDPOINTSOURCE147F4C5896AEDFEFA0ED54AFSLOT179END

The chosen index is a member of the existing canonical family. The source representative, Taylor unit, kernel, action and trace are not chosen again after selecting this index.

For the packet consisting exactly of the quartet in (2), ENDPOINTSOURCE147F4C5896AEDFEFA0ED54AFSLOT180END and ENDPOINTSOURCE147F4C5896AEDFEFA0ED54AFSLOT181END. The earlier exterior bound holds at every admitted index, including the one selected by (39). Since ENDPOINTSOURCE147F4C5896AEDFEFA0ED54AFSLOT182END is increasing on ENDPOINTSOURCE147F4C5896AEDFEFA0ED54AFSLOT183END, (4) follows. Moreover ENDPOINTSOURCE147F4C5896AEDFEFA0ED54AFSLOT184END for every fixed ENDPOINTSOURCE147F4C5896AEDFEFA0ED54AFSLOT185END, which proves the stated threshold ENDPOINTSOURCE147F4C5896AEDFEFA0ED54AFSLOT186END.

A proved upper bound with ENDPOINTSOURCE147F4C5896AEDFEFA0ED54AFSLOT187END would contradict (4). Such an upper bound is \textbf{not} an output of this note. The previous sufficient target ENDPOINTSOURCE147F4C5896AEDFEFA0ED54AFSLOT188END remains sufficient, but (4) specifies the larger quantitative threshold against which a future upper estimate can be compared. No repeated cost-free exterior operation is used: the multiplicity correction in the new formalization is retained.

\subsection{8. The remaining four volumes retain their finite jet maps}\label{endpoint-arithmetic_original--the-remaining-four-volumes-retain-their-finite-jet-maps}

The previous source/relation determinant identity remains

ENDPOINTSOURCE147F4C5896AEDFEFA0ED54AFSLOT189END

In the parallel full-jet transfer coordinates (PR \#24), let ENDPOINTSOURCE147F4C5896AEDFEFA0ED54AFSLOT190END, let ENDPOINTSOURCE147F4C5896AEDFEFA0ED54AFSLOT191END be its ENDPOINTSOURCE147F4C5896AEDFEFA0ED54AFSLOT192END by ENDPOINTSOURCE147F4C5896AEDFEFA0ED54AFSLOT193END raw confluent-jet matrix, and let ENDPOINTSOURCE147F4C5896AEDFEFA0ED54AFSLOT194END be its raw-derivative Vandermonde matrix, with every factorial retained. Its established identity is

ENDPOINTSOURCE147F4C5896AEDFEFA0ED54AFSLOT195END

For the window ENDPOINTSOURCE147F4C5896AEDFEFA0ED54AFSLOT196END, the four transfer indices are exactly ENDPOINTSOURCE147F4C5896AEDFEFA0ED54AFSLOT197END. Put ENDPOINTSOURCE147F4C5896AEDFEFA0ED54AFSLOT198END. Then

ENDPOINTSOURCE147F4C5896AEDFEFA0ED54AFSLOT199END

The common Vandermonde cancels in this explicit ratio, not inside any underlying map. The transfer matrices and their inverses still need their actual arithmetic control; the fixed width alone gives no such bound.

The empty packet has ENDPOINTSOURCE147F4C5896AEDFEFA0ED54AFSLOT200END, ENDPOINTSOURCE147F4C5896AEDFEFA0ED54AFSLOT201END as an empty determinant, and no positive-rank inverse or quartet estimate. The analytic inequalities (23)--(37) still apply to its theta seed, but are not statements about a nonexistent finite packet.

\subsection{9. The volume window is also two explicit retained-boundary costs}\label{endpoint-arithmetic_original--the-volume-window-is-also-two-explicit-retained-boundary-costs}

There is a second useful representation of the same remaining quantity, with no reference-density change. Fix ENDPOINTSOURCE147F4C5896AEDFEFA0ED54AFSLOT202END, both admissible. In the original monic polynomial basis put

ENDPOINTSOURCE147F4C5896AEDFEFA0ED54AFSLOT203END

The kernel update is the actual sum ENDPOINTSOURCE147F4C5896AEDFEFA0ED54AFSLOT204END, where ENDPOINTSOURCE147F4C5896AEDFEFA0ED54AFSLOT205END. The matrix determinant lemma gives

ENDPOINTSOURCE147F4C5896AEDFEFA0ED54AFSLOT206END

This determinant has positive real value. Its factors are ENDPOINTSOURCE147F4C5896AEDFEFA0ED54AFSLOT207END, where ENDPOINTSOURCE147F4C5896AEDFEFA0ED54AFSLOT208END are the generalized eigenvalues of the Hermitian pair ENDPOINTSOURCE147F4C5896AEDFEFA0ED54AFSLOT209END. No non-Hermitian matrix is asserted positive in an unspecified Euclidean metric.

Let ENDPOINTSOURCE147F4C5896AEDFEFA0ED54AFSLOT210END have columns ENDPOINTSOURCE147F4C5896AEDFEFA0ED54AFSLOT211END. The original boundary-layer isomorphism is

ENDPOINTSOURCE147F4C5896AEDFEFA0ED54AFSLOT212END

Its jets vanish; its highest polynomial coefficients give injectivity. Any new boundary can be reduced by these highest coefficients to an old boundary, and orthogonality then gives surjectivity. These statements prove the displayed map with the same original relation spaces, not only a dimension count. The representative update is

ENDPOINTSOURCE147F4C5896AEDFEFA0ED54AFSLOT213END

Every correction in (46) is an original theta boundary. Its split quotient image is ENDPOINTSOURCE147F4C5896AEDFEFA0ED54AFSLOT214END at the receiving fibre, and the input relation remains available in (45). Its cochain primitives are the same fixed-order divisions and tensor signs as in the previous cyclic source construction.

Apply (44) to ENDPOINTSOURCE147F4C5896AEDFEFA0ED54AFSLOT215END and ENDPOINTSOURCE147F4C5896AEDFEFA0ED54AFSLOT216END. This proves

ENDPOINTSOURCE147F4C5896AEDFEFA0ED54AFSLOT217END

In particular

ENDPOINTSOURCE147F4C5896AEDFEFA0ED54AFSLOT218END

where each ENDPOINTSOURCE147F4C5896AEDFEFA0ED54AFSLOT219END has its own indicated source and target from (43). When this trace bound is too large, (47) retains the full logarithmic cost rather than replacing it by its linear upper bound. The analytic norm estimate in (1) does not bound (47); it supplies the other factor in (39).

\subsection{10. Deligne reading and what has actually been transferred}\label{endpoint-arithmetic_original--deligne-reading-and-what-has-actually-been-transferred}

The six supplied S20 parts were reassembled in this turn. Their individual chunk hashes matched, and the assembled archive has 216,580,466 bytes and SHA-256 \texttt{e2005bf31e1fcf362765f73c315c855e0f504f24f34d3a0ab188049da522b7c8}. This verifies the supplied bytes, not every transcription claim in them.

I read the TeX for Definition 1.3.5, Lemmas 3.2.10--3.2.12, the tensor argument 3.2.13, and Corollary 3.3.6 in the nested source files. I compared those exact passages with the supplied top-level French/English frozen page records at printed pages 158, 203, and 206. No PDF text extraction or OCR was used. The source review records a genuine inherited transcription fault: old \texttt{fr\_seg2.tex} writes a weight bound ENDPOINTSOURCE147F4C5896AEDFEFA0ED54AFSLOT220END at 3.2.10, whereas the controlling page record says ENDPOINTSOURCE147F4C5896AEDFEFA0ED54AFSLOT221END; only the latter, with integrality, implies ENDPOINTSOURCE147F4C5896AEDFEFA0ED54AFSLOT222END. The old file also corrupts the explicit squared-eigenvalue step in 3.2.13. Those old files remain unmodified; they are not used as authority for these comparisons. No claim of a full independent audit of Weil II is made.

Deligne's determinant/rank operation and the upper/dual-lower architecture remain attached to the programme's specified exterior and residue maps. Here the additional estimate is a characteristic-zero analytic estimate for the actual source, (1), followed by the endpoint selection (39). This is not a statement that the new base satisfies Deligne's finite-field hypotheses. The remaining volume estimate has not been inferred from purity, positivity of a Gram, or the mere existence of the support.

The new result is therefore a completed analytic component of the current route: its two endpoint monic norms have the required ENDPOINTSOURCE147F4C5896AEDFEFA0ED54AFSLOT223END scale on the actual quartet window. The other session's four-volume theorem now leaves the explicitly specified relation-volume cost as the remaining quantitative input on that window.

\subsection{11. Proof and computation status}\label{endpoint-arithmetic_original--proof-and-computation-status}

Sections 3--7 contain the new written analytic argument. Sections 8--9 contain explicit finite comparisons and a handoff for checking them using the existing quotient and matrix library. The standard complex-analysis, gamma and Legendre facts used in the argument are proved to the extent needed or referenced below. No general novelty claim is made for these classical ingredients.

The accompanying checker tests finite exact identities: coordinate factors, Legendre extremal norms, the local algebra in the interval argument, convolution moments on declared Laplace models, source quotients, block updates, raw confluent-jet endpoint factors, the volume telescope, and retained support cases. These tests do not certify the analytic inequalities by sampling them. There is no new Lean certificate and no interval-certified value of ENDPOINTSOURCE147F4C5896AEDFEFA0ED54AFSLOT224END.

This continuation makes no remote changes. The integration files are add-only, and the parallel formalization branches and the reader's existing sources are left untouched.

\subsection{References and pinned inputs}\label{endpoint-arithmetic_original--references-and-pinned-inputs}

\begin{itemize}
\tightlist
\item
  Deligne, P. \emph{La conjecture de Weil. II}. Publications Mathématiques de l'IHÉS \textbf{52} (1980), 137--252. Supplied S20 source and current page-local records; the selected reading and discrepancies are itemized in \texttt{SOURCE\_REVIEW.md}.
\item
  NIST Digital Library of Mathematical Functions: §§25.4 (zeta functional equation), 5.11 (gamma estimates), 18.3 and 18.5 (Legendre polynomials). Public source pages consulted: https://dlmf.nist.gov/25.4, https://dlmf.nist.gov/5.11, https://dlmf.nist.gov/18.3, https://dlmf.nist.gov/18.5.
\item
  Original programme: the delivered \emph{Tau Gamma Convolution Descent}, \emph{Tau Toda Volume Control}, \emph{Tau Cyclic Sum Control}, and \emph{Tau Exterior Trace Amplification} notes. Their definitions of the source, action, unit and canonical quotient are retained; this is not a complete replay of all their upstream analytic claims.
\item
  Parallel endpoint criterion: PR \#23, \texttt{c720f40530eed2f5969dbabe94dfd3fddc0f507f}, \texttt{workbenches/tau-toda-volume-formal/ENDPOINT\_CRITERION.md}.
\item
  Parallel confluent transfer: PR \#24, \texttt{dfcbba5cbf7fec8c9301fe242c13e741d762013e}, \texttt{workbenches/tau-confluent-transfer/RESEARCH\_NOTE.md}.
\item
  Inspected main: \texttt{b32ca2128e0deb0eec6eafb860776a5d6a28dcbb}, after its merge of the original-Gram exterior/trace contribution.
\end{itemize}
